% GENERATED FILE. Edit canonical lessons, then run npm run generate:books.
% canonical-source-hash: fnv1a64:df2bd30c199a5e19
% canonical-lessons: ES-C688-oceano, ES-C688-ola, ES-C688-marea, ES-C688-bahia, ES-C688-flotar

\chapter{El mar}
\label{ch:es-el-mar}
\hlchaptermodality{688}

\begin{chapteropening}
\textbf{By the end of this chapter:} \emph{I can talk about the sea: the ocean, waves, tides, bays and floating in the water.}

Complete the last lesson of chapter 688: I can talk about the sea: the ocean, waves, tides, bays and floating in the water.
\end{chapteropening}

\section[el océano]{el océano — the ocean}

\label{lesson:ES-C688-oceano}



\begin{quote}
Before the new one: say the Spanish for to clamber, then the Spanish for to explore.
\end{quote}

\subsection*{You'll want to know: el océano}
\textbf{el océano} — \textquotedblleft{}the ocean\textquotedblright{}.

\begin{quote}
El océano está muy frío.
\end{quote}

\begin{quote}
\emph{The ocean is very cold.}
\end{quote}

\emph{El océano} is bigger than \emph{el mar}.

\begin{cousinweb}
\emph{océano} is a learned borrowing from Latin \emph{ōceanus}, from Greek \emph{Ōkeanós} (EtymDB), the great river that the Greeks thought ran round the world.
\end{cousinweb}

\begin{grammarlens}[title={What you've built}]
One more word for telling, explaining and describing at B1.
\end{grammarlens}

\subsection*{Your turn}
\begin{itemize}
\raggedright
  \item \emph{Say it:} \emph{el océano}
  \item \emph{Say it:} \emph{el océano}, once more
  \item \emph{Say it:} a sentence of your own with \emph{el océano}
  \item \emph{Recall:} say the Spanish for to clamber, then the Spanish for to explore, then say \emph{el océano} again
\end{itemize}

\begin{tcolorbox}[breakable,colback=teal!4,colframe=teal!35!black,title={Before you move on}]
What does el océano mean? (\textquotedblleft{}The ocean\textquotedblright{}.) Say it once more.
\end{tcolorbox}
\section[la ola]{la ola — the wave}

\label{lesson:ES-C688-ola}



\begin{quote}
Before the new one: say the Spanish for to explore, then the Spanish for the ocean.
\end{quote}

\subsection*{You'll want to know: la ola}
\textbf{la ola} — \textquotedblleft{}the wave\textquotedblright{}.

\begin{quote}
Hoy hay olas muy grandes.
\end{quote}

\begin{quote}
\emph{The waves are very big today.}
\end{quote}

\emph{Una ola de calor} — a heat wave.

\begin{cousinweb}
The origin of \emph{ola} is not settled; several sources have been proposed, none of them sure.
\end{cousinweb}

\begin{grammarlens}[title={What you've built}]
One more word for telling, explaining and describing at B1.
\end{grammarlens}

\subsection*{Your turn}
\begin{itemize}
\raggedright
  \item \emph{Say it:} \emph{la ola}
  \item \emph{Say it:} \emph{la ola}, once more
  \item \emph{Say it:} a sentence of your own with \emph{la ola}
  \item \emph{Recall:} say the Spanish for to explore, then the Spanish for the ocean, then say \emph{la ola} again
\end{itemize}

\begin{tcolorbox}[breakable,colback=teal!4,colframe=teal!35!black,title={Before you move on}]
What does la ola mean? (\textquotedblleft{}The wave\textquotedblright{}.) Say it once more.
\end{tcolorbox}
\section[la marea]{la marea — the tide}

\label{lesson:ES-C688-marea}



\begin{quote}
Before the new one: say the Spanish for the ocean, then the Spanish for the wave.
\end{quote}

\subsection*{You'll want to know: la marea}
\textbf{la marea} — \textquotedblleft{}the tide\textquotedblright{}.

\begin{quote}
Con la marea baja se puede pasar a la isla.
\end{quote}

\begin{quote}
\emph{At low tide you can walk across to the island.}
\end{quote}

\emph{Marea alta} / \emph{baja} — high / low tide.

\begin{cousinweb}
\emph{marea} was borrowed from French \emph{marée} \textquotedblleft{}tide\textquotedblright{}, built on \emph{mer} \textquotedblleft{}sea\textquotedblright{}, from Latin \emph{mare} (EtymDB).
\end{cousinweb}

\begin{grammarlens}[title={What you've built}]
One more word for telling, explaining and describing at B1.
\end{grammarlens}

\subsection*{Your turn}
\begin{itemize}
\raggedright
  \item \emph{Say it:} \emph{la marea}
  \item \emph{Say it:} \emph{la marea}, once more
  \item \emph{Say it:} a sentence of your own with \emph{la marea}
  \item \emph{Recall:} say the Spanish for the ocean, then the Spanish for the wave, then say \emph{la marea} again
\end{itemize}

\begin{tcolorbox}[breakable,colback=teal!4,colframe=teal!35!black,title={Before you move on}]
What does la marea mean? (\textquotedblleft{}The tide\textquotedblright{}.) Say it once more.
\end{tcolorbox}
\section[la bahía]{la bahía — the bay}

\label{lesson:ES-C688-bahia}



\begin{quote}
Before the new one: say the Spanish for the wave, then the Spanish for the tide.
\end{quote}

\subsection*{You'll want to know: la bahía}
\textbf{la bahía} — \textquotedblleft{}the bay\textquotedblright{}.

\begin{quote}
El hotel está en una bahía tranquila.
\end{quote}

\begin{quote}
\emph{The hotel is on a quiet bay.}
\end{quote}

A small bay is \emph{una cala}.

\begin{cousinweb}
\emph{bahía} is usually traced to French \emph{baie} \textquotedblleft{}bay\textquotedblright{}; where the French word comes from is not settled.
\end{cousinweb}

\begin{grammarlens}[title={What you've built}]
One more word for telling, explaining and describing at B1.
\end{grammarlens}

\subsection*{Your turn}
\begin{itemize}
\raggedright
  \item \emph{Say it:} \emph{la bahía}
  \item \emph{Say it:} \emph{la bahía}, once more
  \item \emph{Say it:} a sentence of your own with \emph{la bahía}
  \item \emph{Recall:} say the Spanish for the wave, then the Spanish for the tide, then say \emph{la bahía} again
\end{itemize}

\begin{tcolorbox}[breakable,colback=teal!4,colframe=teal!35!black,title={Before you move on}]
What does la bahía mean? (\textquotedblleft{}The bay\textquotedblright{}.) Say it once more.
\end{tcolorbox}
\section[flotar]{flotar — to float}

\label{lesson:ES-C688-flotar}



\begin{quote}
Before the new one: say the Spanish for the tide, then the Spanish for the bay.
\end{quote}

\subsection*{You'll want to know: flotar}
\textbf{flotar} — \textquotedblleft{}to float\textquotedblright{}.

\begin{quote}
En este mar flotas sin esfuerzo.
\end{quote}

\begin{quote}
\emph{In this sea you float easily.}
\end{quote}

\emph{El flotador} is a rubber ring.

\begin{cousinweb}
\emph{flotar} was borrowed from French \emph{flotter} \textquotedblleft{}to float\textquotedblright{} (EtymDB), a word of Germanic origin, a cousin of English \emph{float}.
\end{cousinweb}

\begin{grammarlens}[title={What you've built}]
One more word for telling, explaining and describing at B1.
\end{grammarlens}

\subsection*{Your turn}
\begin{itemize}
\raggedright
  \item \emph{Say it:} \emph{flotar}
  \item \emph{Say it:} \emph{flotar}, once more
  \item \emph{Say it:} a sentence of your own with \emph{flotar}
  \item \emph{Recall:} say the Spanish for the tide, then the Spanish for the bay, then say \emph{flotar} again
\end{itemize}

\begin{tcolorbox}[breakable,colback=teal!4,colframe=teal!35!black,title={Before you move on}]
What does flotar mean? (\textquotedblleft{}To float\textquotedblright{}.) Say it once more.
\end{tcolorbox}
